\section{Reference Interfaces and Data Contracts}
\label{app:interfaces}

The following interfaces are normative at the semantic level. Names may change before a stable release, but an implementation must preserve the ownership and resource contracts.

\subsection{Core value types}

The descriptors below are immutable values shared by planning, execution, storage, diagnostics, and telemetry; they carry no implicit ownership.

\begin{lstlisting}[language=Python,caption={Core immutable descriptors.}]
from __future__ import annotations
from dataclasses import dataclass
from enum import Enum
from typing import Mapping, Sequence, Tuple

class NumericMode(str, Enum):
    REFERENCE = "reference"
    DETERMINISTIC = "deterministic"
    STABLE = "stable"
    FAST = "fast"
    APPROXIMATE = "approximate"

class TensorClass(str, Enum):
    PARAMETER = "parameter"
    ACTIVATION = "activation"
    ACCUMULATOR = "accumulator"
    KV_STATE = "kv_state"
    OUTPUT = "output"
    GRADIENT = "gradient"
    OPTIMIZER = "optimizer"

@dataclass(frozen=True, slots=True)
class Tile:
    origin: Tuple[int, ...]
    extent: Tuple[int, ...]

    def validate(self, shape: Sequence[int]) -> None:
        if len(self.origin) != len(shape) or len(self.extent) != len(shape):
            raise ValueError("tile rank mismatch")
        for o, e, n in zip(self.origin, self.extent, shape, strict=True):
            if o < 0 or e < 0 or o > n or e > n - o:
                raise ValueError("tile outside logical tensor")

@dataclass(frozen=True, slots=True)
class ResourceVector:
    bytes_by_arena: Mapping[str, int]
    credits_by_queue: Mapping[str, int]

@dataclass(frozen=True, slots=True)
class VirtualTensorSpec:
    tensor_id: str
    storage_id: str
    shape: Tuple[int, ...]
    logical_dtype: str
    storage_dtype: str
    strides: Tuple[int, ...]
    storage_offset_elements: int
    tensor_class: TensorClass
    immutable: bool
    codec: str
    version: int

def checked_nbytes(shape: Sequence[int], itemsize: int, limit: int) -> int:
    if itemsize <= 0:
        raise ValueError("invalid item size")
    value = itemsize
    for dim in shape:
        if dim < 0 or (dim and value > limit // dim):
            raise OverflowError("tensor byte size exceeds configured limit")
        value *= dim
    return value
\end{lstlisting}

The actual implementation uses fixed-width checked integers in the native core. Python descriptors are validated before crossing the ABI.

\subsection{Backend protocol}

The backend boundary exposes bounded primitives and isolates vendor-specific streams, events, copies, kernels, and runtime reservations.

\begin{lstlisting}[language=Python,caption={Backend contract.}]
from typing import Protocol, Iterable, Optional

@dataclass(frozen=True, slots=True)
class KernelCapability:
    op: str
    implementation_id: str
    dtypes: Tuple[str, ...]
    deterministic: bool
    min_alignment: int
    shape_multiples: Tuple[int, ...]
    max_workspace_bytes: int

@dataclass(frozen=True, slots=True)
class BackendCapabilities:
    backend_id: str
    devices: Tuple[str, ...]
    arena_alignment: int
    async_copy: bool
    copy_compute_overlap: bool
    direct_storage: bool
    kernels: Tuple[KernelCapability, ...]

class Event(Protocol):
    @property
    def terminal(self) -> bool: ...
    def wait(self, timeout_s: Optional[float] = None) -> None: ...
    def error(self) -> Optional[BaseException]: ...

class BufferLease(Protocol):
    @property
    def arena_id(self) -> str: ...
    @property
    def size_bytes(self) -> int: ...
    @property
    def epoch(self) -> int: ...
    def bind_terminal_event(self, event: Event) -> None: ...

class Backend(Protocol):
    def capabilities(self) -> BackendCapabilities: ...
    def reserve_runtime(self, config: "BackendConfig") -> "BackendReservation": ...
    def copy_async(
        self, src: BufferLease, dst: BufferLease, *, stream: str
    ) -> Event: ...
    def launch(
        self,
        implementation_id: str,
        inputs: Sequence[BufferLease],
        outputs: Sequence[BufferLease],
        workspace: Optional[BufferLease],
        attributes: Mapping[str, object],
        *,
        stream: str,
        wait_for: Sequence[Event],
    ) -> Event: ...
    def quiesce(self) -> None: ...
\end{lstlisting}

\code{launch} may not allocate material workspace outside the supplied lease. Any fixed context allocation occurs in \code{reserve\_runtime} and is charged to the backend reserve.

\subsection{Storage protocol}

The storage boundary is range-oriented, cancellation-aware, integrity-checking, and transactional so that no API requires whole-tensor materialization.

\begin{lstlisting}[language=Python,caption={Range-oriented storage contract.}]
@dataclass(frozen=True, slots=True)
class ByteRange:
    object_id: str
    offset: int
    length: int
    checksum: str | None

@dataclass(frozen=True, slots=True)
class IORequest:
    ranges: Tuple[ByteRange, ...]
    priority: int
    deadline_class: str
    direct: bool
    cancellation_id: str

class TensorStore(Protocol):
    def capabilities(self) -> "StoreCapabilities": ...
    def read_async(self, req: IORequest, dst: BufferLease) -> Event: ...
    def write_async(
        self, req: IORequest, src: BufferLease, *, transaction_id: str
    ) -> Event: ...
    def publish(self, transaction_id: str, root_record: bytes) -> None: ...
    def abort(self, transaction_id: str) -> None: ...
    def close(self) -> None: ...
\end{lstlisting}

The store validates that the sum of range lengths equals or maps safely into the destination, and it reports short IO as an error. A direct path may impose stricter alignment; the planner sees those constraints.

\subsection{Operator plugin protocol}

An operator plugin owns semantic inference, minimum-working-set calculation, legal plan enumeration, reference behavior, and optional backward lowering.

\begin{lstlisting}[language=Python,caption={Semantic operator and plan-provider contract.}]
@dataclass(frozen=True, slots=True)
class PlanContext:
    capabilities: BackendCapabilities
    budgets: ResourceVector
    numeric_mode: NumericMode
    request_bounds: Mapping[str, int]
    layouts: Mapping[str, object]
    calibration_id: str

@dataclass(frozen=True, slots=True)
class OperatorPlan:
    implementation_id: str
    resource_bound: ResourceVector
    tile_parameters: Mapping[str, int]
    task_template: bytes
    numerical_contract: Mapping[str, object]
    estimated_cost_ns: int

class OperatorPlugin(Protocol):
    op_name: str
    schema_version: int

    def infer(
        self, inputs: Sequence[VirtualTensorSpec], attrs: Mapping[str, object]
    ) -> Sequence[VirtualTensorSpec]: ...

    def minimum_working_set(
        self, inputs: Sequence[VirtualTensorSpec], attrs: Mapping[str, object],
        ctx: PlanContext,
    ) -> ResourceVector: ...

    def enumerate_plans(
        self, inputs: Sequence[VirtualTensorSpec], attrs: Mapping[str, object],
        ctx: PlanContext,
    ) -> Iterable[OperatorPlan]: ...

    def reference(
        self, logical_inputs: Sequence[object], attrs: Mapping[str, object]
    ) -> Sequence[object]: ...

    def backward_plugin(self) -> "OperatorPlugin | None": ...
\end{lstlisting}

A conformance runner checks that every emitted plan's observed allocations fit \code{resource\_bound}, including tail tiles and error paths.

\subsection{Engine API}

The public API separates model opening and preflight reservation from request execution, making the bounded-residency contract observable before the first token.

\begin{lstlisting}[language=Python,caption={Public model/session API.}]
@dataclass(frozen=True, slots=True)
class RequestEnvelope:
    max_prompt_tokens: int
    max_new_tokens: int
    max_batch: int = 1
    max_beams: int = 1
    max_speculative_tokens: int = 0

@dataclass(frozen=True, slots=True)
class PreflightResult:
    accepted: bool
    plan_id: str | None
    reservation_id: str | None
    memory_bound_bytes: Mapping[str, int]
    spill_bound_bytes: int
    estimated_physical_bytes: Mapping[str, int]
    fallback_chain: Tuple[str, ...]
    rejection: "AmsError | None"

class ModelHandle(Protocol):
    def preflight(
        self, envelope: RequestEnvelope, policy: "ExecutionPolicy"
    ) -> PreflightResult: ...
    def open_session(
        self, preflight: PreflightResult, request: "GenerationRequest"
    ) -> "Session": ...
    def close(self) -> None: ...

class Session(Protocol):
    def next_token(self) -> "TokenResult": ...
    def stream(self) -> "Iterable[TokenResult]": ...
    def cancel(self) -> None: ...
    def checkpoint(self, uri: str) -> None: ...
    def close(self) -> None: ...
\end{lstlisting}

\code{preflight} is required. A convenience \code{generate} may call it internally, but it returns the accepted plan summary and never executes under an unvalidated implicit budget.

\subsection{Structured errors}

Stable structured errors preserve machine-actionable category, phase, resource evidence, fallback history, and remediation details across language bindings.

\begin{lstlisting}[language=Python,caption={Stable error payload.}]
@dataclass(frozen=True, slots=True)
class AmsError:
    code: str
    message: str
    retryable: bool
    phase: str
    operator: str | None
    graph_location: str | None
    logical_tile: Tile | None
    plan_id: str | None
    cause_chain: Tuple[str, ...]
    diagnostics: Mapping[str, object]
\end{lstlisting}

Diagnostics may contain sizes, capacities, file IDs, and backend codes, but never tensor contents, secrets, or raw prompts by default. Stable top-level codes are versioned; backend-specific details are nested.

\subsection{Native ABI requirements}

The C ABI passes opaque handles and plain fixed-layout structs with explicit size/version fields. It must provide:

\begin{itemize}
  \item create/destroy backend and store instances;
  \item capability enumeration;
  \item reserve/release arenas;
  \item lease buffer and bind event;
  \item submit copy/kernel/read/write;
  \item poll/wait/cancel/quiesce;
  \item return structured status without C++ exceptions crossing the boundary.
\end{itemize}

Every function documents thread safety, ownership transfer, reentrancy, callback context, and whether it may block. ABI conformance is fuzzed with invalid handles, sizes, versions, and callback races.
